% !PRINTED 48
\weberpage{48}

\begin{center}
\S.\,12.

\gesperrt{Divisors of an Abelian Group}.

\gesperrt{Reciprocal Groups}.
\end{center}

A divisor $T$ of an \gesperrt{Abel}ian group $S$ is, as already remarked above, always a normal divisor, and hence by \S.\,4 there is a group complementary to $T$,
\[
S/T,
\]
whose elements are the subsidiary groups of $T$, namely
\weq{1}{T,\ TA',\ TA''\ldots,}
where $A', A'' \ldots$ denote certain elements of $S$. The number of the elements (1), that is, the degree of the group $S/T$, is, if $t$ is the degree of $T$, equal to the index of the divisor $T$ of $S$, hence equal to
\[
\frac{n}{t} = j.
\]

But this group $S/T$ is itself again an \gesperrt{Abel}ian one; for we have
\[
TA'\,TA'' = TA'A'',\ \ TA''\,TA' = TA''A',
\]
so that the two are equal to each other.

We have now to determine the characters of the group $S/T$. We denote the elements of this group by
\weq{2}{T,\ T_1,\ T_2 \ldots T_{j-1},}
and one of their characters by $\xi(T_i)$.

From this function $\xi(T)$ we can now derive a function $\xi(A)$ of the elements of $S$ by setting
\weq{3}{\xi(A_i) = \xi(T_i)}
whenever $A_i$ is any element occurring in the subsidiary group $T_i$. For the elements of the group $T$ itself we then have $\xi(A) = 1$.

But this function $\xi(A_i)$ is contained among the characters of $S$. For if $A_i$ occurs in $T_i$ and $A_k$ in $T_k$, then $A_iA_k$ occurs in $T_iT_k$, and consequently
\weq{4}{\xi(A_i)\xi(A_k) = \xi(T_i)\xi(T_k) = \xi(T_iT_k) = \xi(A_iA_k).}
But this, by \S.\,11,(4), is the definition of a character of $S$.
